\section{Introduction}
\label{sec:intro}

Binary64 arithmetic is fast and widely available, but its precision is not
enough for every numerical computation \citep{Higham2002}.  Arbitrary-precision
libraries such as MPFR address this limitation by allowing each variable to
have an independently chosen precision \citep{Fousse2007}.  When roughly twice
the precision of binary64 is sufficient, double-double arithmetic offers a
more specialized alternative.  It represents a number as the unevaluated sum
of two binary64 values and performs arithmetic on the pair.  This approach has
been used to improve reproducibility in parallel climate models
and to provide additional precision in quantum electrodynamic calculations
\citep{HeDing2001,Aoyama2015}.

The apparent simplicity of the representation hides an important hazard.
\citet{Li2002} described double-double arithmetic as ``an `attractive
nuisance,' like an unfenced backyard swimming pool.''  A pair of binary64
values does not form an ordinary wider IEEE~754 format, and a double-double
operation is not obtained by computing an exact result and rounding it once.
Instead, a double-double operation is implemented as a sequence of rounded
binary64 computations.  Different algorithms can implement the same
double-double operation, and the error bound depends on the chosen algorithm.
A separate argument must also show that the returned pair is a valid
double-double.  The error bound and validity guarantee hold only under
conditions on the operands and intermediate results.

We present a verified C library whose public C interface enforces these
conditions.  The library provides three error-free transformations and six
double-double operations.  On normal return, each double-double operation produces 
a valid double-double within its proved error bound, while each error-free 
transformation guarantees exact reconstruction.

\subsection{Double-Word Arithmetic and Its Numerical Guarantees}
\label{sec:intro-double-word}

We state the numerical analysis for double-word arithmetic over a generic
underlying floating-point format.  Double-double arithmetic is the binary64
instance of double-word arithmetic.  A double word is a floating-point
expansion with two components \citep{Dekker1971,Priest1991}.  Let
$\mathbb{F}$ denote the underlying format, and let $\operatorname{RN}$ denote
round-to-nearest, ties-to-even. Following
\citet[Definition~1.4]{Joldes2017}, we use the following definition.

\begin{definition}[Double word]
\label[definition]{def:double-word}
A double word $x$ over $\mathbb{F}$ is the unevaluated sum
\begin{equation}
  x=x_h+x_\ell
\end{equation}
of two values $x_h,x_\ell\in\mathbb{F}$ such that
\begin{equation}
  x_h=\operatorname{RN}(x).
  \label{eq:dw-representation}
\end{equation}
\end{definition}

Thus, rounding the represented value $x$ recovers the high component $x_h$,
while the low component is the exact residual $x_\ell=x-x_h$.  For an exact
real result $r$, a double-word algorithm should return $(z_h,z_\ell)$
satisfying two distinct properties.  Writing $\gamma$ for an error bound
derived by the analysis, the accuracy property is
\begin{equation}
  \left|(z_h+z_\ell)-r\right|\leq\gamma,
  \label{eq:intro-accuracy-property}
\end{equation}
and the validity property is that the returned components form a double word
$z$ according to \cref{def:double-word}.  The validity property supplies the
representation precondition required when the result is used by a later
double-word operation.

The numerical results on which we build
\citep{Joldes2017,Rideau2022} establish both properties for double-word
addition, multiplication, and division when underflow and overflow do not
occur.  \citet{Rideau2022} mechanize these results using Flocq's FLX format
\citep{Boldo2011}, whose unbounded exponent range represents this assumption.

Native IEEE~754 arithmetic requires both ends of the exponent range to be
considered.  Near the upper end, an intermediate computation can overflow even when its
operands are finite, and the returned pair may consequently contain an
infinity or NaN.  Neither an infinity nor a NaN represents a real value, so
such a pair can satisfy neither the accuracy property nor the valid
double-word property.  Both properties
therefore require the returned components to be finite.  Near the lower end,
the fixed spacing of subnormal values can invalidate exact relations used by
the basic building blocks of double-word algorithms and prevent rounding
errors from being bounded solely relative to the exact result.

Our analysis determines how these effects change the guarantee for each
operation.  Addition retains the relative bounds of the earlier analysis.
Multiplication and division instead satisfy mixed bounds: a relative term
accounts for ordinary rounding error, while an absolute term accounts for
errors introduced near the smallest representable values.  For every
operation, the resulting theorem gives sufficient conditions under which the
returned pair is a valid double word and its represented value satisfies the
stated error bound.

Conditions governing internal computations must be established before
execution.  Result finiteness can be established in either of two
ways: stronger conditions on the inputs can imply it before execution, or the
returned components can be checked afterward.  Making the numerical results
available through ordinary C calls therefore requires runtime checks that
establish the appropriate conditions before or after the kernel call.

\subsection{This Work: A Verified C Library for Double-Double Arithmetic}

For each of the six double-double operations, the library provides two public
wrappers around the same arithmetic kernel.  The default wrapper performs a
weak-precondition check before execution, invokes the kernel, and then checks
that both returned components are finite.  The other wrapper performs a
strong-precondition check that establishes result finiteness before execution,
so no result check is needed afterward.  The terms \emph{weak} and
\emph{strong} describe the precondition checks, not the wrappers.  Both
wrappers provide the same contract on normal return: the represented result
satisfies the proved error bound, and the returned pair is a valid double word
over binary64---that is, a valid double-double result.

For each operation, we define in Rocq a functional model over a generic
radix-two format and prove its numerical guarantee using Flocq.  We then use
VST to prove that the C kernel implements the binary64 instance of this model
and that the runtime checks establish the conditions required by its
numerical theorem.  Together, these proofs establish the contracts of the
public C interface.

We evaluate whether the resulting checks are both effective and practical.
Across tests spanning the binary64 exponent range, every double-double result
returned normally was valid and within its proved error bound; the largest
observed error was \RQOneMaximumBoundRatio{} times the corresponding bound.
For calls stopped by a runtime check, the experiment ran the corresponding QD
operation on the same inputs and measured whether its result satisfied the
corresponding numerical criterion.  Finally,
the performance measurements separate the cost of the runtime checks from the
cost of the arithmetic kernels.  \Cref{sec:evaluation} gives the experimental
methods and results.

\subsection{Contributions}

In summary, this paper makes four contributions:

\begin{itemize}
  \item \textbf{Runtime checks for double-double arithmetic
        (\Cref{sec:overview}).}
        We implement public wrappers whose runtime checks establish the
        sufficient conditions required by these numerical results.

  \item \textbf{Numerical guarantees over IEEE~754
        (\Cref{sec:numerics}).}
        We extend prior mechanized analysis to the finite exponent range of
        IEEE~754, proving error bounds and valid double-word results for six
        operations and guarantees of exact reconstruction for three
        error-free transformations.

  \item \textbf{A verified C implementation of double-double arithmetic
        (\Cref{sec:mechanization}).}
        Using VST, we connect the Rocq functional models and numerical
        theorems to the binary64 C kernels, runtime checks, and contracts of
        the public C interface.

  \item \textbf{Empirical evaluation
        (\Cref{sec:evaluation}).}
        We evaluate the tightness of the proved bounds, the behavior of QD on
        inputs associated with calls for which a runtime check failed, and the
        cost of the public wrappers relative to the arithmetic kernels and QD.
\end{itemize}

The accompanying artifact contains the C sources, Rocq and VST proofs, and
evaluation scripts at \url{ARTIFACT-URL}, revision
\texttt{ARTIFACT-COMMIT}.
